% \documentclass{article}
\documentclass[border=0mm, convert={density=300, quality=90, outext=.png}]{standalone}
% ===== Loading all packages =====
% === input, languages, layout ===
\usepackage[utf8]{inputenc}
\usepackage[T1]{fontenc}
\usepackage[ngerman]{babel}
\usepackage{textcomp}
\usepackage{url}
\usepackage{printlen}
\usepackage{lmodern}                % Latin Modern, bessere Fonts

% === mathematics, numbers, units ===
\usepackage{amsmath}
\usepackage{amssymb}
\usepackage{siunitx}
\usepackage{calc}                   % enables the calculation of distances etc.

% === figures, plots, drawing, ... ===
\usepackage[cmyk]{xcolor}
\usepackage{graphicx}
\usepackage{tikz}
\usepackage{pgfplots}
\usepackage{pgfplotstable}

% === other stuff ===
\usepackage{multirow}               % Table entries reaching over more than one line
\usepackage[
    layout = inline,
    final]{fixme}                  	% comments
\usepackage{ifthen}

% === setspace ===
\linespread{1.15}

% === tikz ===
\usetikzlibrary{
    fit,
    positioning,
    shapes,
    calc,
    pgfplots.groupplots,
    plotmarks,
    backgrounds,
    matrix,
    fadings,
    external,
    decorations.markings,
    patterns,
    arrows.meta,
    intersections}

% === xcolor ===
% Color set "strong" of TU Dortmund university
\definecolor{tuGreen}{cmyk}{0.57, 0, 1, 0}
\definecolor{tuOrange}{cmyk}{0, 0.55, 1, 0}
\definecolor{tuBlue}{cmyk}{1, 0, 0.36, 0}
\definecolor{tuViolette}{cmyk}{0.25, 1, 0, 0}
\definecolor{tuYellow}{cmyk}{0.05, 0.08, 1, 0}
\definecolor{tuBrown}{cmyk}{0.05, 0.69, 1, 0}
\definecolor{red}{rgb}{1, 0, 0}
\definecolor{blue}{rgb}{0, 0, 1}

% === pgfplots ===
% Definition of cycle lists
\pgfplotscreateplotcyclelist{TUDefault}{%
    {plain plot set = {none}{tuGreen}},%
    {plain plot set = {none}{tuOrange}},%
    {plain plot set = {none}{tuBlue}},%
    {plain plot set = {none}{tuViolette}},%
    {plain plot set = {none}{tuYellow}},%
    {plain plot set = {none}{tuBrown}}%
}

\pgfplotsset{
    compat = newest,
    scale only axis,
    width = 0.85\textwidth,
    height = 0.25\textheight,
    every axis legend/.append style = {nodes={right}},
    xmajorgrids,
    ymajorgrids,
    legend style = {font = \tiny, at = {(0.5,1.025)}, anchor = south, legend columns = 6},
    % Custom plot style (https://tex.stackexchange.com/questions/193589/pgfplots-how-can-i-define-one-cycle-list-for-all-graphs-or-the-least-amount-p/193596)
    plain plot set/.style 2 args = {
        #2,
        mark = #1,
        every mark/.append style = {fill = #2}
    },
    % Axis labels according to DIN 461 (based on https://tex.stackexchange.com/questions/391332/get-the-auto-generated-tick-distance-of-a-pgfplot-din-461/391430#391430)
    % Argument 1: Axis description
    % Argument 2: Axis unit
    din xlabel/.style 2 args={
        xticklabel style={
            name=xlabel\ticknum,
            append after command=\pgfextra{\xdef\lastxticknum{\ticknum}}
        },
        after end axis/.append code={
            \pgfmathparse{int(\lastxticknum-1)}
            \path (xlabel\lastxticknum.base) -- (xlabel\pgfmathresult.base) node (xAxUnit) [midway, anchor=base] {#2};
            \node[anchor = east] (xAxDescr) at ($(xlabel\lastxticknum.base) + (-5mm,-3mm)$) {#1};
            \draw[->] (xAxDescr.east) -- ($(xlabel\lastxticknum.base) + (0,-3mm)$);
        }
    },
    din ylabel/.style 2 args={
        yticklabel style={
            name=ylabel\ticknum,
            append after command=\pgfextra{\xdef\lastyticknum{\ticknum}}
        },
        after end axis/.append code={
            \pgfmathparse{int(\lastyticknum-1)}
            \path (ylabel\lastyticknum.base) -- (ylabel\pgfmathresult.base) node (yAxUnit) [midway, anchor=base] {#2};
            \node[anchor = east, rotate = 90] (yAxDescr) at ($(ylabel\lastyticknum.west) + (-3mm,-5mm)$) {#1};
            \draw[->] (yAxDescr.east) -- ($(ylabel\lastyticknum.west) + (-3mm,0)$);
        }
    },
    din zlabel/.style 2 args={
        zticklabel style={
            name=zlabel\ticknum,
            append after command=\pgfextra{\xdef\lastzticknum{\ticknum}}
        },
        after end axis/.append code={
            \pgfmathparse{int(\lastzticknum-1)}
            \path (zlabel\lastzticknum.base) -- (zlabel\pgfmathresult.base) node (zAxUnit) [midway, anchor=base] {#2};
            \node[anchor = east, rotate = 90] (zAxDescr) at ($(zlabel\lastzticknum.west) + (-3mm,-5mm)$) {#1};
            \draw[->] (zAxDescr.east) -- ($(zlabel\lastzticknum.west) + (-3mm,0)$);
        }
    }
}

\usepgfplotslibrary{groupplots}
\pgfdeclarelayer{background}
\pgfdeclarelayer{foreground}
\pgfsetlayers{background,main,foreground}
 
% === SIunitx ===
\sisetup{
    load-configurations = abbreviations,            % load units with abbreviations
    range-phrase = \ldots,
    list-final-separator = { \translate{und} },
    list-pair-separator = { \translate{und} },
    output-decimal-marker = {,},
    group-minimum-digits = 3,
    group-separator = {.},
    exponent-product = \cdot,
    per-mode = symbol-or-fraction
}
\DeclareSIUnit{\Wp}{Wp}
\DeclareSIUnit{\kWp}{kWp}
\DeclareSIUnit{\kWh}{kWh}
\DeclareSIUnit{\GWh}{GWh}
\DeclareSIUnit{\TWh}{TWh}
\DeclareSIUnit{\pu}{p.u.}
\DeclareSIUnit{\VA}{VA}
\DeclareSIUnit{\VAr}{VAr}

% Command for printing the complex j
\newcommand{\cj}{\text{j}}

\begin{document}
    \begin{tikzpicture}
    % === Anschlüsse ===
        \node[circle, draw = tuGreen, fill = tuGreen, inner sep = 0, minimum height = 1.5mm] (port_a) at (0,0){};
        \node[circle, draw = tuOrange, fill = tuGreen, inner sep = 0, minimum height = 1.5mm] (port_b) at (15mm,0){};
        \node[circle, draw = tuOrange, fill = tuGreen, inner sep = 0, minimum height = 1.5mm] (port_c) at (30mm,0){};
        \node[circle, draw = tuOrange, fill = tuGreen, inner sep = 0, minimum height = 1.5mm] (port_d) at (45mm,0){};
        \node[circle, draw = tuOrange, fill = tuOrange, inner sep = 0, minimum height = 1.5mm] (port_e) at (72.5mm,0){};

        \draw (port_a.west) -- ++(-3mm, -1mm) -- ++(0mm, -2mm) edge[densely dotted] ++(0mm, -1.8mm);
        \draw (port_a.west) -- ++(-3mm, 1mm) -- ++(0mm, 2mm) edge[densely dotted] ++(0mm, 1.8mm);
        \draw (port_a.east) -- ++(4.5mm,0) -- ++(4.5mm, 1.5mm) ++(0,-1.5mm) -- (port_b.west);
        \draw (port_b.east) -- ++(4.5mm,0) -- ++(4.5mm, 1.5mm) ++(0,-1.5mm) -- (port_c.west);
        \draw (port_c.east) -- ++(4.5mm,0) -- ++(4.5mm, 1.5mm) ++(0.,-1.5mm) -- (port_d.west);

        \node[circle, draw = black, inner sep = 0mm, minimum height = 5mm, anchor = west] (winding_a) at (55.0mm,0){};
        \node[circle, draw = black, inner sep = 0mm, minimum height = 5mm, anchor = west] (winding_b) at (57.5mm,0){};

        \draw (port_d.east) -- (winding_a.west) (winding_b.east) -- (port_e.west);

        \node[tuGreen, anchor = north, inner sep = 2mm, text width = 10mm, align = center] at (port_a.south) {A \\ \SI{110}{\kV} \\ 1};
        \node[tuGreen, anchor = north, inner sep = 2mm, text width = 10mm, align = center] at (port_b.south) {B \\ \SI{110}{\kV} \\ 1};
        \node[tuGreen, anchor = north, inner sep = 2mm, text width = 10mm, align = center] at (port_c.south) {C \\ \SI{110}{\kV} \\ 1};
        \node[tuGreen, anchor = north, inner sep = 2mm, text width = 10mm, align = center] at (port_d.south) {D \\ \SI{110}{\kV} \\ 1};
        \node[tuOrange, anchor = north, inner sep = 2mm, text width = 10mm, align = center] at (port_e.south) {E \\ \SI{10}{\kV} \\ 2};

            % Corner points of SubGridContainerA
            \node[anchor=south west, font=\footnotesize] (SGC_A_SW) at (-0.5, 0.75) {SubGridContainer A}; % South West corner
            \node[anchor=south east] (SGC_A_SE) at (5.25, 1.0) {}; % South East corner
            \node[anchor=south west] (SGC_A_NW) at (-0.5, -1.5) {}; % North West corner
            \node[anchor=south east] (SGC_A_NE) at (5.25, -1.5) {}; % North East corner

            % Corner points of SubGridContainerB
            \node[anchor=south west] (SGC_B_NW) at (4.0, 0.75) {}; % North West corner
            \node[anchor=south east, font=\footnotesize] (SGC_B_NE) at (8.25, 0.75) {SubGridContainer B}; % North East corner
            \node[anchor=north west] (SGC_B_SW) at (4.0, -2.0) {}; % South West corner
            \node[anchor=south east] (SGC_B_SE) at (8.25, -2.0) {}; % South East corner

            % red rectangle
            \draw [red] (SGC_A_SW.north west) -- (SGC_A_SE.north east) -- (SGC_A_NE.south east) -- (SGC_A_NW.south west) -- cycle;
            % blue rectangle
            \draw [blue] (SGC_B_SW.north west) -- (SGC_B_SE.south east) -- (SGC_B_NE.south east) -- (SGC_B_NW.south west) -- cycle;
	    \end{tikzpicture}
\end{document}